\newpage

\section{মাযীদ}

\begin{center}
  {\bfseries ছুলাসী মাযীদ, হামযায় ওয়াসল}
  \enspace
  {\arabicfont \textarabic{بهمزة وصل}}
\end{center}

\vspace{0.1cm}

\noindent নয় বাব। শুরুতে হামযায় ওয়াসল। চিহ্ন বাড়তি অক্ষর।

\vspace{0.15cm}

\begingroup
\setlength{\arrayrulewidth}{0.5pt}
\setlength{\tabcolsep}{3pt}
\renewcommand{\arraystretch}{1.42}
\small

\noindent\begin{tabularx}{\textwidth}{%
    |>{\hsize=0.26\hsize\centering\arraybackslash}X%
    |>{\hsize=1.28\hsize\centering\arraybackslash}X%
    |>{\hsize=1.46\hsize\raggedright\arraybackslash}X|}
  \hline
  \rowcolor{headergreen}
  {\color{white}\textbf{বাব}} &
  {\color{white}\textbf{মাজি · মুজারে}} &
  {\color{white}\textbf{চিহ্ন}} \\
  \hline
  \brow{lightgreen}{১}{%
    {\arabicfont \textarabic{اِجْتَنَبَ يَجْتَنِبُ}}\newline দূরে থাকা}{%
    {\arabicfont \textarabic{اِفْتِعَال}}\newline
    ফার আগে হামযা, ফা-আইনে তা}
  \brow{white}{২}{%
    {\arabicfont \textarabic{اِسْتَنْصَرَ يَسْتَنْصِرُ}}\newline সাহায্য চাওয়া}{%
    {\arabicfont \textarabic{اِسْتِفْعَال}}\newline
    হামযা, সীন, তা। লাযিম ও মুতাআদ্দি}
  \brow{lightgreen}{৩}{%
    {\arabicfont \textarabic{اِنْفَطَرَ يَنْفَطِرُ}}\newline বিদীর্ণ হওয়া}{%
    {\arabicfont \textarabic{اِنْفِعَال}}\newline
    হামযা ও নুন। লাযিম}
  \brow{white}{৪}{%
    {\arabicfont \textarabic{اِحْمَرَّ يَحْمَرُّ}}\newline রক্তিম হওয়া}{%
    {\arabicfont \textarabic{اِفْعِلَال}}\newline
    হামযা, লামে শাদ্দা। লাযিম}
  \brow{lightgreen}{৫}{%
    {\arabicfont \textarabic{اِدْهَامَّ يَدْهَامُّ}}\newline খুব কালো হওয়া}{%
    {\arabicfont \textarabic{اِفْعِيلَال}}\newline
    হামযা, আলিফ, লামে শাদ্দা। লাযিম}
  \brow{white}{৬}{%
    {\arabicfont \textarabic{اِخْشَوْشَنَ يَخْشَوْشِنُ}}\newline অত্যন্ত শক্ত হওয়া}{%
    {\arabicfont \textarabic{اِفْعِيعَال}}\newline
    হামযা, আইন দুবার, মাঝে ওয়াও। লাযিম}
  \brow{lightgreen}{৭}{%
    {\arabicfont \textarabic{اِجْلَوَّذَ يَجْلَوِّذُ}}\newline উট দৌড়ে চলা}{%
    {\arabicfont \textarabic{اِفْعِوَّال}}\newline
    হামযা, আইনের পর শাদ্দা-ওয়াও। লাযিম}
  \brow{white}{৮}{%
    {\arabicfont \textarabic{اِثَّاقَلَ يَثَّاقَلُ}}\newline ভারি হওয়া}{%
    {\arabicfont \textarabic{اِفَّاعَلَ}}\newline
    হামযা, ফায় শাদ্দা, মাঝে আলিফ। লাযিম}
  \brow{lightgreen}{৯}{%
    {\arabicfont \textarabic{اِطَّهَّرَ يَطَّهَّرُ}}\newline পবিত্র হওয়া}{%
    {\arabicfont \textarabic{اِفَّعَّلَ}}\newline
    হামযা, ফা ও আইনে শাদ্দা। লাযিম}
\end{tabularx}
\endgroup

\vspace{0.35cm}

\begin{center}
  {\bfseries হামযায় ওয়াসল নেই}
  \enspace
  {\arabicfont \textarabic{بغير همزة وصل}}
\end{center}

\vspace{0.1cm}

\noindent পাঁচ বাব। প্রথম বাবে হামযায় কাতঈ, বাকি বাবে তা ও আলিফ।
{\arabicfont \textarabic{تَفَعَّلَ}} ও {\arabicfont \textarabic{تَفَاعَلَ}}-এর শুরুতে দুই তা মিললে এক তা বাদ যায়।

\vspace{0.15cm}

\begingroup
\setlength{\arrayrulewidth}{0.5pt}
\setlength{\tabcolsep}{3pt}
\renewcommand{\arraystretch}{1.42}
\small

\noindent\begin{tabularx}{\textwidth}{%
    |>{\hsize=0.26\hsize\centering\arraybackslash}X%
    |>{\hsize=1.28\hsize\centering\arraybackslash}X%
    |>{\hsize=1.46\hsize\raggedright\arraybackslash}X|}
  \hline
  \rowcolor{headergreen}
  {\color{white}\textbf{বাব}} &
  {\color{white}\textbf{মাজি · মুজারে}} &
  {\color{white}\textbf{চিহ্ন}} \\
  \hline
  \brow{lightgreen}{১}{%
    {\arabicfont \textarabic{أَكْرَمَ يُكْرِمُ}}\newline সম্মান করা}{%
    {\arabicfont \textarabic{إِفْعَال}}\newline
    হামযায় কাতঈ। মুজারে হামযা যায়, হরফে পেশ}
  \brow{white}{২}{%
    {\arabicfont \textarabic{صَرَّفَ يُصَرِّفُ}}\newline রূপান্তর করা}{%
    {\arabicfont \textarabic{تَفْعِيل}}\newline
    আইনে শাদ্দা। মুজারার হরফে পেশ}
  \brow{lightgreen}{৩}{%
    {\arabicfont \textarabic{تَقَبَّلَ يَتَقَبَّلُ}}\newline গ্রহণ করা}{%
    {\arabicfont \textarabic{تَفَعُّل}}\newline
    ফার আগে তা, আইনে শাদ্দা}
  \brow{white}{৪}{%
    {\arabicfont \textarabic{قَاتَلَ يُقَاتِلُ}}\newline পরস্পর লড়াই}{%
    {\arabicfont \textarabic{مُفَاعَلَة}}\newline
    ফার পর আলিফ। মুজারার হরফে পেশ}
  \brow{lightgreen}{৫}{%
    {\arabicfont \textarabic{تَقَابَلَ يَتَقَابَلُ}}\newline পরস্পর মুখোমুখি}{%
    {\arabicfont \textarabic{تَفَاعُل}}\newline
    ফার আগে তা, ফার পর আলিফ}
\end{tabularx}
\endgroup

\vspace{0.2cm}

\noindent পরে আলাদা ছকে আছে
{\arabicfont \textarabic{تَفْعِيل،}}
{\arabicfont \textarabic{تَفَعُّل،}}
{\arabicfont \textarabic{اِفْتِعَال}}
ও {\arabicfont \textarabic{اِسْتِفْعَال}}।

\vspace{0.4cm}

\begin{center}
  {\bfseries রুবায়ী}
  \enspace
  {\arabicfont \textarabic{رباعي}}
\end{center}

\vspace{0.1cm}

\noindent মুজারাদ এক বাব, মুনশাইব তিন। হামযাওয়ালা দুই বাব লাযিম।

\vspace{0.15cm}

\begingroup
\setlength{\arrayrulewidth}{0.5pt}
\setlength{\tabcolsep}{3pt}
\renewcommand{\arraystretch}{1.42}
\small

\noindent\begin{tabularx}{\textwidth}{%
    |>{\hsize=0.26\hsize\centering\arraybackslash}X%
    |>{\hsize=1.28\hsize\centering\arraybackslash}X%
    |>{\hsize=1.46\hsize\raggedright\arraybackslash}X|}
  \hline
  \rowcolor{headergreen}
  \multicolumn{3}{|c|}{\color{white}\textbf{মুজারাদ}} \\
  \hline
  \brow{lightgreen}{১}{%
    {\arabicfont \textarabic{بَعْثَرَ يُبَعْثِرُ}}\newline উত্তেজিত করা}{%
    {\arabicfont \textarabic{فَعْلَلَة}}\newline
    মূল চার। মুজারার হরফে পেশ, শেষের আগে যের}
  \rowcolor{headergreen}
  \multicolumn{3}{|c|}{\color{white}\textbf{মুনশাইব, হামযা নেই}} \\
  \hline
  \brow{white}{১}{%
    {\arabicfont \textarabic{تَسَرْبَلَ يَتَسَرْبَلُ}}\newline জামা পরিধান}{%
    {\arabicfont \textarabic{تَفَعْلُل}}\newline
    ফার আগে তা। লাযিম}
  \rowcolor{headergreen}
  \multicolumn{3}{|c|}{\color{white}\textbf{মুনশাইব, হামযায় ওয়াসল}} \\
  \hline
  \brow{lightgreen}{১}{%
    {\arabicfont \textarabic{اِبْرَنْشَقَ يَبْرَنْشِقُ}}\newline খুশি হওয়া}{%
    {\arabicfont \textarabic{اِفْعِنْلَال}}\newline
    হামযা, আইন-লামে নুন। লাযিম}
  \brow{white}{২}{%
    {\arabicfont \textarabic{اِقْشَعَرَّ يَقْشَعِرُّ}}\newline রোমাঞ্চিত হওয়া}{%
    {\arabicfont \textarabic{اِفْعِلَالّ}}\newline
    হামযা, শেষ লামে শাদ্দা। লাযিম}
\end{tabularx}
\endgroup
